%% ma: L12-B06-0082
%% lop: 12
%% chuong: 2
%% bai: 6
%% dang: 12.6.22
%% loai: TN4
%% muc: VD
%% dap_an: "A"
%% nguon: "CĐ2.6.pdf giáo viên gửi 27/10/2026 (M07, Lê Minh Khôi – Nguyễn Anh Khôi, Chương II Chủ đề 6 Ứng dụng vào các bài toán thực tế), Câu 25"
%% kien_thuc: "Bài 6 (tổng vectơ, bình phương vô hướng), lớp 10 (phương trình bậc hai)"
%% trang_thai: cho-duyet
%% ly_do: "Dạng mới 12.6.22 đề xuất, chờ giáo viên duyệt"
%% ghi_chu: "Đề gốc dùng chung giả thiết cho Câu 25–26; đã chép vào từng câu"
\begin{ex}
Trong không gian, cho bốn sợi dây buộc với nhau tại điểm $I$. Điểm $I$ cân bằng dưới tác dụng của bốn lực dọc theo bốn sợi dây lần lượt là $\vec F_1$, $\vec F_2$, $\vec F_3$, $\vec F_4$. Biết rằng $|\vec F_1|=180$ N, $|\vec F_3|=220$ N, góc tạo bởi hai lực $\left(\vec F_2,\vec F_4\right)=120^\circ$, góc tạo bởi hai lực $\left(\vec F_1,\vec F_3\right)=50^\circ$. Đơn vị lực trong bài toán tính theo Newton. Biết độ lớn của lực $\vec F_2$ bằng $180$ N. Tính độ lớn của lực $\vec F_4$ (làm tròn kết quả đến hàng đơn vị).
\choice{\True $418$ N}{$315$ N}{$363$ N}{$400$ N}
\loigiai{Cân bằng: $\vec F_1+\vec F_3=-\left(\vec F_2+\vec F_4\right)$ nên $\left|\vec F_1+\vec F_3\right|^2=\left|\vec F_2+\vec F_4\right|^2$.

$180^2+220^2+2\cdot180\cdot220\cos50^\circ\approx131711$; với $x=|\vec F_4|$: $180^2+x^2-180x=131711$.

$x^2-180x-99311=0\Rightarrow x\approx418$ N.}
\end{ex}
